\documentclass[a4paper,12pt]{article}
\usepackage[utf8]{inputenc}
\usepackage{amsmath}
\usepackage{geometry}
\geometry{margin=1in}

\begin{document}
\pagestyle{empty}
% Encabezado
\begin{center}
    \textbf{\huge Recuperatorio} \\[0.5cm]
    \textbf{\large Escuela N° 4-124 Escuela Técnica "Reynaldo Merín"} \\[0.5cm]
    \textbf{Curso: Electrotecnia II 4°1°} \\[0.5cm]
    \textbf{Nombre y Apellido: \underline{\hspace{10cm}}} \\[0.5cm]
    \textbf{Fecha: \underline{\hspace{3cm}}} \\[1cm]
\end{center}

\noindent \textbf{Instrucciones:} Resuelva los siguientes problemas. Cada problema tiene un valor de 1.43 puntos. Escriba sus respuestas de manera clara y ordenada, mostrando todos los pasos necesarios para llegar a la solución.

\vspace{0.5cm}

% Problemas
\begin{enumerate}
    \item \textbf{Problema 1: (1.43 puntos)} \\
    Un conductor rectilíneo de longitud \(L = 0.5 \, \text{m}\) está ubicado en un campo magnético uniforme de \(B = 0.2 \, \text{T}\). Si una corriente de \(I = 10 \, \text{A}\) fluye a través del conductor, determine la fuerza ejercida sobre el conductor.
    
    \item \textbf{Problema 2: (1.43 puntos)} \\
    Según la ley de Faraday-Lenz, ¿cuál es la fuerza electromotriz inducida en una espira circular de radio \(r = 0.1 \, \text{m}\) cuando el campo magnético a través de la espira cambia uniformemente de \(B = 0.5 \, \text{T}\) a \(B = 0 \, \text{T}\) en 2 segundos?
    
    \item \textbf{Problema 3: (1.43 puntos)} \\
    Un conductor rectilíneo de longitud \(l = 1 \, \text{m}\) se mueve a una velocidad de \(v = 3 \, \text{m/s}\) perpendicularmente a un campo magnético de \(B = 0.4 \, \text{T}\). Calcule la fuerza electromotriz inducida en el conductor.
    
    \item \textbf{Problema 4: (1.43 puntos)} \\
    Un inductor de inductancia \(L_1 = 5 \, \text{H}\) está en serie con otro inductor de inductancia \(L_2 = 3 \, \text{H}\). ¿Cuál es la inductancia equivalente del sistema?
    
    \item \textbf{Problema 5: (1.43 puntos)} \\
    Calcule la autoinducción L de una bobina de \(N = 100\) vueltas, longitud \(l = 0.2 \, \text{m}\), y área de sección transversal \(A = 0.01 \, \text{m}^2\), sabiendo que la permeabilidad del núcleo es \(\mu = 4 \pi \times 10^{-7} \, \text{H/m}\).
    
    \item \textbf{Problema 6: (1.43 puntos)} \\
    Determine f.e.m. en una bobina que tiene una autoinducción de \(L = 2 \, \text{H}\) cuando se pasa una corriente de \(I = 5 \, \text{A}\) a \(I = 0 \, \text{A}\) en 1 segundo.
    
    \item \textbf{Problema 7: (1.43 puntos)} \\
    Dos capacitores \(C_1 = 20 \, \mu\text{F}\) y \(C_2 = 30 \, \mu\text{F}\) están conectados en serie. Determine la capacidad equivalente.
\end{enumerate}

\vfill
\begin{center}
    \textbf{¡Buena suerte!}
\end{center}

\end{document}
